% NOT FOR THE EXAM

\section{The Gaussian Model}

The Gaussian model is a simple yet powerful statistical model that assumes that the data follows a normal distribution. It is widely used in various fields, including physics, engineering, and machine learning, due to its mathematical tractability and the central limit theorem.

It is a very limited case of application of RG methods, but it is a good starting point to understand the basic concepts and techniques of RG. In this section, we will explore the Gaussian model in detail, including its formulation, properties, and applications.

The gaussian model is defined by keeping only the quadratic term in the LG partition function:
\[
    \mathcal{Z} = \int \mathcal{D}m(\mathbf{x}) \, \exp {-\int \mathrm{d}^D \mathbf{x} \, \left[ \cdots \right]},
\]
and it is a bit unphysical since it is only defined for \(t \geq 0\), but it is a good starting point to understand the basic concepts and techniques of RG. We would not even need RG to solve this model, but it is a good example to illustrate the RG procedure.

We will provide a \textbf{momentum space RG solution} for this model.

    [...]

\subsection{Exact Solution}

The Gaussian model is exactly solvable, and we can find the exact solution by performing a Fourier transform of the field \(m(\mathbf{x})\). The Fourier transform allows us to express the field in terms of its momentum components, which simplifies the analysis of the partition function. The Fourier transform of the field \(m(\mathbf{x})\) is defined as:
\[
    \tilde{m}(q) = \int \mathrm{d}^D \mathbf{x} \, m(\mathbf{x}) e^{i \mathbf{q} \cdot \mathbf{x}}, \iff m(\mathbf{x}) = \sum_{\mathbf{q}} e^{-i \mathbf{q} \cdot \mathbf{x}} \tilde{m}(\mathbf{q}) \sim \int {\mathrm{d}^D \mathbf{q}}{(2\pi)} \, e^{-i \mathbf{q} \cdot \mathbf{x}} \tilde{m}(\mathbf{q}),
\]
where the sum is over the discrete momenta allowed by the finite size of the system, and the integral is over the continuous momenta in the thermodynamic limit. The partition function can be rewritten plugging the Fourier transform of \(m(\mathbf{x})\) into the original expression:
\[
    \mathcal{Z} = e^{\frac{V\hbar^2}{2t}}\dots
\]
If we take the logarithm, the product becomes a sum which we can express as an integral over the Brillouin zone: since the singularities of the integrand are at \(q=0\), we can replace the Brillouin zone with a sphere of radius \(\Lambda\) (the UV cutoff): ...

And in the limit for \(t \to 0\) ...

Therefore we can identify
\[
    \begin{dcases}
        \alpha_+ = 2 - \frac{D}{2}, \\
        \Delta = \frac{1}{2} + \frac{D}{4}.
    \end{dcases}
\]

[...]

\begin{enumerate}
    \item \textbf{Coarse Grain}

          [...]

    \item \textbf{Rescale}

          [...]

    \item \textbf{Renormalize}

          [...]
\end{enumerate}

\subsubsection{Relevant and Irrelevant Operators}

[...]

\begin{remark}[Dangerous Irrelevant Operators]

\end{remark}

\subsubsection{Infinitesimal Scale Transformations}

[...]